\documentclass[12pt]{article}
\usepackage[a4paper, margin=1in]{geometry}
\usepackage{amsmath, amssymb}
\usepackage{enumitem}

\title{Lecture 26: Stochastic Processes -- Part 1}
\author{}
\date{}

\begin{document}

\maketitle

\section*{Topics}
SSS, WSS, Ergodic Process

\section{Basics of Stochastic Processes (Revision)}

\subsection{Concept of Waveform Collections and Realizations}

\begin{itemize}
    \item A stochastic process is viewed as a \textbf{collection of waveforms}.
    \item Each waveform corresponds to a \textbf{realization}.
    \item Each realization is associated with an outcome of the experiment.
\end{itemize}

\subsubsection*{Logical Understanding:}

\begin{itemize}
    \item Fix outcome $\rightarrow$ get one waveform $\rightarrow$ called a realization
    \item Fix time $\rightarrow$ get a random variable
\end{itemize}

Thus:

\begin{itemize}
    \item \textbf{Along time $\rightarrow$ waveform}
    \item \textbf{At fixed time $\rightarrow$ random variable}
\end{itemize}

\subsection{Realizations and Random Variables}

\begin{itemize}
    \item A stochastic process connects:
    \begin{itemize}
        \item Time domain behavior (waveforms)
        \item Probabilistic behavior (random variables)
    \end{itemize}
\end{itemize}

\subsubsection*{Key Interpretation:}

\begin{itemize}
    \item A process can be seen as:
    \begin{itemize}
        \item A \textbf{family of random variables indexed by time}
        \item OR a \textbf{collection of sample functions}
    \end{itemize}
\end{itemize}

\subsection{Characterization of Stochastic Process}

The process is described using:

\begin{itemize}
    \item First-order density function
    \item Second-order density function
    \item N-th order density function
\end{itemize}

\subsubsection*{Logical Flow:}

\begin{itemize}
    \item To fully describe a process:
    $\rightarrow$ Need joint distributions at multiple time instants
    \item Increasing order $\rightarrow$ more complete description
\end{itemize}

\section{Statistical Moments}

\subsection{Mean of Stochastic Process}

\begin{itemize}
    \item Mean is defined as:
    \begin{itemize}
        \item Expected value at each time instant
    \end{itemize}
\end{itemize}

\subsubsection*{Interpretation:}

\begin{itemize}
    \item Provides \textbf{average behavior over realizations}
    \item Depends on time in general
\end{itemize}

\subsection{Autocorrelation Function}

\begin{itemize}
    \item Measures similarity between values at two time instants
\end{itemize}

\subsubsection*{Logical Role:}

\begin{itemize}
    \item Captures \textbf{dependency structure}
    \item Helps understand:
    \begin{itemize}
        \item Memory of process
        \item Predictability
    \end{itemize}
\end{itemize}

\subsection{Properties of Autocorrelation}

From slides (structure only):

\begin{itemize}
    \item Symmetry-related properties
    \item Dependence on time indices
    \item Relation to covariance
\end{itemize}

\subsection{Covariance Function}

\begin{itemize}
    \item Defined using mean and autocorrelation
\end{itemize}

\subsubsection*{Logical Relation:}

\begin{itemize}
    \item Covariance = measure of \textbf{fluctuation around mean}
\end{itemize}

\section{Positive Definite Functions and Matrices}

\subsection{Positive Definite Function}

\begin{itemize}
    \item Introduced via autocorrelation properties
\end{itemize}

\subsubsection*{Key Idea:}

\begin{itemize}
    \item Autocorrelation must satisfy \textbf{positive definiteness}
\end{itemize}

\subsection{Matrix Expansion}

\begin{itemize}
    \item Autocorrelation values arranged in matrix form
\end{itemize}

\subsubsection*{Logical Flow:}

\begin{enumerate}
    \item Construct matrix using autocorrelation values
    \item Apply quadratic form
    \item Show non-negativity
\end{enumerate}

\subsection{Quadratic Form Analogy}

\begin{itemize}
    \item Used to interpret positive definiteness
\end{itemize}

\subsubsection*{Extension:}

\begin{itemize}
    \item From 2D $\rightarrow$ 3D interpretation (as per slides)
\end{itemize}

\subsection{Positive Definite Matrices}

\begin{itemize}
    \item Definition based on quadratic form
\end{itemize}

\subsubsection*{Methods Mentioned:}

\begin{itemize}
    \item Elimination method
    \item Matrix-based verification
\end{itemize}

\subsection{Example (Positive Definite Matrix)}

\begin{itemize}
    \item Slides include:
    \begin{itemize}
        \item Example 1
        \item Continued steps
    \end{itemize}
\end{itemize}

\subsubsection*{Important Note:}

\begin{itemize}
    \item Since derivation is not present in text, only structure is retained:
    \begin{itemize}
        \item Matrix formation
        \item Stepwise checking
        \item Final conclusion
    \end{itemize}
\end{itemize}

\section{Stationary Stochastic Processes}

\subsection{Concept of Stationarity (Time Invariance)}

\begin{itemize}
    \item Stationarity implies \textbf{statistical properties do not change with time}
\end{itemize}

\subsection{Types of Stationarity}

\begin{itemize}
    \item Strict Sense Stationary (SSS)
    \item Wide Sense Stationary (WSS)
\end{itemize}

\subsection{Strict Sense Stationary (SSS)}

\subsubsection*{Definition (from structure):}

\begin{itemize}
    \item All joint distributions are invariant under time shift
\end{itemize}

\subsubsection*{Logical Meaning:}

\begin{itemize}
    \item Complete statistical behavior remains unchanged
\end{itemize}

\subsection{Wide Sense Stationary (WSS)}

\subsubsection*{Definition (from structure):}

\begin{itemize}
    \item Only:
    \begin{itemize}
        \item Mean is constant
        \item Autocorrelation depends on time difference
    \end{itemize}
\end{itemize}

\subsection{WSS vs SSS}

\subsubsection*{Implied Comparison:}

\begin{itemize}
    \item SSS $\rightarrow$ stronger condition
    \item WSS $\rightarrow$ weaker condition
\end{itemize}

\subsection{Gaussian Process: WSS $\Rightarrow$ SSS?}

\begin{itemize}
    \item Slides explicitly raise this question
\end{itemize}

\subsubsection*{Logical Insight:}

\begin{itemize}
    \item For Gaussian processes:
    \begin{itemize}
        \item WSS may imply SSS (as discussed conceptually in slides)
    \end{itemize}
\end{itemize}

\subsection{Example Problems}

Slides include:

\begin{itemize}
    \item Example $\rightarrow$ Solution $\rightarrow$ continuation
\end{itemize}

\subsubsection*{Structure of Solution:}

\begin{enumerate}
    \item Identify process properties
    \item Compute mean
    \item Compute autocorrelation
    \item Check stationarity conditions
\end{enumerate}

(No explicit formulas provided in text, so not expanded)

\subsection{HW Example: White Noise Process}

\begin{itemize}
    \item Mentioned as a standard example under WSS
\end{itemize}

\section{Ergodic Process}

\subsection{Concept}

\begin{itemize}
    \item Introduced after stationarity
\end{itemize}

\subsubsection*{Key Idea:}

\begin{itemize}
    \item Relates \textbf{time averages} and \textbf{ensemble averages}
\end{itemize}

\subsection{Time Average of Random Process}

Slides indicate stepwise development:

\begin{enumerate}
    \item Define time averaging
    \item Compare with statistical average
    \item Establish condition for equality
\end{enumerate}

\subsection{Key Logical Flow:}

\begin{itemize}
    \item If: Time average = Ensemble average
    $\rightarrow$ Process is ergodic
\end{itemize}

\section{Example 4 (Detailed in Slides)}

\subsection*{Structure of Solution:}

\subsubsection*{Step 1: Mean Calculation}

\begin{itemize}
    \item Multiple slides devoted to:
    \begin{itemize}
        \item Mean derivation
        \item Continuation steps
    \end{itemize}
\end{itemize}

\subsubsection*{Step 2: Autocorrelation}

\begin{itemize}
    \item Step-by-step expansion
    \item Multiple continuation slides
\end{itemize}

\subsubsection*{Step 3: Final Result}

\begin{itemize}
    \item Concluding expression given
\end{itemize}

\textbf{Note:} Since explicit expressions are not visible in text, only logical progression is preserved.

\section{Mathematical Tools for Studying Random Processes}

Slides mention tools but do not expand in text.

\section{Final Revision Concepts}

\subsection{Generalized Representation}

\begin{itemize}
    \item Process expressed in generalized mathematical form
\end{itemize}

\subsection{Key Takeaways}

\begin{itemize}
    \item Stochastic process = collection of random variables
    \item Mean and autocorrelation define second-order behavior
    \item Positive definiteness is essential for validity
    \item Stationarity simplifies analysis
    \item Ergodicity connects theory with observation
\end{itemize}

\section*{Exam Preparation Strategy (Based on This Lecture)}

Focus on:

\begin{enumerate}
    \item Definitions:
    \begin{itemize}
        \item SSS, WSS, Ergodic
    \end{itemize}
    \item Conceptual differences:
    \begin{itemize}
        \item WSS vs SSS
    \end{itemize}
    \item Logical conditions:
    \begin{itemize}
        \item Time invariance
        \item Dependence on time difference
    \end{itemize}
    \item Properties:
    \begin{itemize}
        \item Autocorrelation
        \item Positive definiteness
    \end{itemize}
    \item Problem-solving flow:
    \begin{itemize}
        \item Mean $\rightarrow$ Autocorrelation $\rightarrow$ Check stationarity
    \end{itemize}
\end{enumerate}

\end{document}